% ============================================================
%  Manual Técnico — Proyecto K-QGMIP
%  Universidad de Caldas
%  Análisis y Diseño de Algoritmos — 2026-1
% ============================================================
\documentclass[11pt,letterpaper]{article}

% ── Paquetes ─────────────────────────────────────────────────────────────────
\usepackage[utf8]{inputenc}
\usepackage[T1]{fontenc}
\usepackage[spanish]{babel}
\usepackage{times}
\usepackage[margin=2.5cm]{geometry}
\usepackage{amsmath,amssymb,amsthm}
\usepackage{mathtools}
\usepackage{algorithm}
\usepackage{algpseudocode}
\usepackage{listings}
\usepackage{xcolor}
\usepackage{graphicx}
\usepackage{float}
\usepackage{booktabs}
\usepackage{longtable}
\usepackage{multirow}
\usepackage{hyperref}
\usepackage{tikz}
\usetikzlibrary{shapes,arrows.meta,positioning,fit,calc,matrix}
\usepackage{pgfplots}
\pgfplotsset{compat=1.18}
\usepackage{caption}
\usepackage{subcaption}
\usepackage{fancyhdr}
\usepackage{setspace}
\usepackage{enumitem}
\usepackage{csquotes}

% ── Colores ───────────────────────────────────────────────────────────────────
\definecolor{codebg}{rgb}{0.95,0.95,0.97}
\definecolor{codeframe}{rgb}{0.7,0.7,0.7}
\definecolor{keycolor}{rgb}{0.0,0.4,0.7}
\definecolor{commentcolor}{rgb}{0.4,0.6,0.4}
\definecolor{stringcolor}{rgb}{0.7,0.3,0.1}

% ── Estilo de listados (código Python) ────────────────────────────────────────
\lstdefinestyle{pycode}{
  language=Python,
  backgroundcolor=\color{codebg},
  frame=single,
  rulecolor=\color{codeframe},
  basicstyle=\small\ttfamily,
  keywordstyle=\color{keycolor}\bfseries,
  commentstyle=\color{commentcolor}\itshape,
  stringstyle=\color{stringcolor},
  numberstyle=\tiny\color{gray},
  numbers=left,
  numbersep=5pt,
  tabsize=4,
  breaklines=true,
  showstringspaces=false,
  captionpos=b,
}
\lstset{style=pycode}

% ── Estilo de pseudocódigo ────────────────────────────────────────────────────
\algrenewcommand\algorithmicrequire{\textbf{Entrada:}}
\algrenewcommand\algorithmicensure{\textbf{Salida:}}
\algnewcommand\algorithmiccomment[1]{\hfill\(\triangleright\) #1}

% ── Cabecera y pie ────────────────────────────────────────────────────────────
\pagestyle{fancy}
\fancyhf{}
\fancyhead[L]{\small Manual Técnico — K-QGMIP}
\fancyhead[R]{\small Universidad de Caldas 2026-1}
\fancyfoot[C]{\thepage}
\renewcommand{\headrulewidth}{0.4pt}

% ── Comandos matemáticos ─────────────────────────────────────────────────────
\newcommand{\EMD}{\mathrm{EMD}}
\newcommand{\MIP}{\mathrm{MIP}}
\newcommand{\kMIP}{k\text{-}\mathrm{MIP}}
\newcommand{\phi}{\varphi}
\newcommand{\Phi}{\Phi}
\newcommand{\tpm}{\mathbf{T}}
\newcommand{\parti}{\mathcal{P}}

% ── Teoremas ──────────────────────────────────────────────────────────────────
\newtheorem{definition}{Definición}[section]
\newtheorem{theorem}{Teorema}[section]
\newtheorem{proposition}{Proposición}[section]
\newtheorem{corollary}{Corolario}[section]

% ============================================================
\begin{document}

% ── Portada ───────────────────────────────────────────────────────────────────
\begin{titlepage}
  \centering
  \vspace*{1cm}
  {\large \textbf{Universidad de Caldas}\\[0.3cm]
   Facultad de Inteligencia Artificial e Ingenierías\\[0.3cm]
   Análisis y Diseño de Algoritmos --- 2026-1}
  
  \vspace{3cm}
  {\Huge \textbf{Manual Técnico}}\\[0.5cm]
  {\LARGE \textbf{Proyecto K-QGMIP}}\\[0.8cm]
  {\large \textit{Extensión a k-particiones de los algoritmos\\
   KGeoMIP y KQNodes para el cálculo del\\
   System Irreducibility Analysis (SIA/IIT)}}
  
  \vspace{3cm}
  {\large Fecha: \today}
  
  \vfill
  \begin{abstract}
    Este manual describe el diseño, implementación y evaluación experimental
    del proyecto \textbf{K-QGMIP}: una extensión de los algoritmos GeoMIP y
    QNodes para encontrar la Minimum Information Partition (MIP) generalizada
    a $k$ grupos en redes con $n$ nodos. Se presentan los fundamentos
    matemáticos del System Irreducibility Analysis (SIA), la arquitectura del
    software, los algoritmos pseudocodificados, el análisis de complejidad
    asintótica y los resultados experimentales sobre matrices de transición
    de hasta $n=22$ nodos.
  \end{abstract}
\end{titlepage}

% ── Tabla de contenidos ───────────────────────────────────────────────────────
\tableofcontents
\newpage

% ============================================================
\section{Resumen Ejecutivo}
\label{sec:resumen}
% ============================================================

\subsection{Problema abordado}

El análisis de irreducibilidad de sistemas (SIA) cuantifica cuánta información
se \emph{pierde} al particionar un sistema de $n$ nodos en $k$ grupos
desconectados.  Formalmente, dado un sistema definido por su
Transition Probability Matrix (TPM) $\tpm \in [0,1]^{2^n \times n}$ y un estado
inicial $s_0 \in \{0,1\}^n$, se busca la $k$-partición
$\parti^* = \{P_1,\ldots,P_k\}$ de los nodos que minimiza la pérdida de
información integrada:
\[
  \phi(\parti) = \EMD\!\bigl(p(S_t \mid s_0),\;
                             \bigotimes_{i=1}^k p(P_i^{t} \mid s_0^{P_i})\bigr),
\]
siendo $\EMD$ la Earth Mover's Distance (distancia de Wasserstein de orden 1)
entre el repertorio efecto del subsistema completo y el producto tensorial de
los repertorios de cada partición~\cite{tononi2016}.

\subsection{Estrategias desarrolladas}

\begin{description}[style=nextline]
  \item[\textbf{KGeoMIP} (\texttt{KPartitionSIA})]
    Extensión de GeoMIP~\cite{geomip2025} a $k$ particiones mediante
    biparticiones greedy encadenadas y clustering jerárquico aglomerativo.
    Complejidad $\mathcal{O}(k \cdot n^2 \cdot 2^n)$.
  \item[\textbf{KQNodes} (\texttt{QNodes} extendido)]
    Extensión heurística de QNodes que aplica $k-1$ cortes greedy usando
    el coste marginal de cada nodo.  Complejidad $\mathcal{O}(k \cdot n^2)$.
  \item[\textbf{GeometricSIA} (base, $k=2$)]
    Programación dinámica sobre el hipercubo de estados; garantiza el MIP óptimo
    para biparticiones en $\mathcal{O}(n \cdot 2^n)$.
  \item[\textbf{BruteForce} (referencia)]
    Enumeración exhaustiva de todas las biparticiones.
    $\mathcal{O}(2^{n-1})$; solo viable para $n \le 6$.
\end{description}

\subsection{Resultados principales}

\begin{itemize}
  \item GeometricSIA reproduce el Cuadro 4.2 del paper GeoMIP con error
        $< 10^{-4}$ para el sistema N3C (validación cruzada con 34 tests).
  \item GeometricSIA y BruteForce coinciden en N3A (5 estados iniciales
        evaluados, diferencia $< 10^{-6}$).
  \item KPartitionSIA devuelve $k$ y $n_\text{nodos}$ correctos en todos los
        casos de prueba.
  \item Suite de 34 tests automatizados: \textbf{100\% de paso}.
\end{itemize}

\subsection{Limitaciones}

\begin{itemize}
  \item La heurística greedy de KPartitionSIA no garantiza
        $\phi(k{=}3) \le \phi(k{=}2)$ en todos los casos (aproximadamente
        33\,\% de los casos en N3--N5).
  \item Para $n > 15$ la memoria requerida ($\mathcal{O}(2^n \cdot n)$) supera
        los 16 GB; se recomienda reducir el subsistema o usar paralelismo
        (ver Apéndice~\ref{app:gpu}).
\end{itemize}

% ============================================================
\section{Fundamentos Teóricos}
\label{sec:teoria}
% ============================================================

\subsection{Integrated Information Theory (IIT) y SIA}

La Integrated Information Theory propuesta por Tononi~\cite{tononi2016} modela
un sistema como un conjunto de $n$ variables aleatorias binarias
$\mathbf{X} = (X_1,\ldots,X_n)$ con distribución conjunta dada por la TPM.
El \emph{repertorio efecto} del subsistema bajo el mecanismo $M \subseteq V$
y el alcance $Z \subseteq V$ desde el estado $s$ se define como:
\[
  p(Z^{t+1} \mid M^t = s^M, \overline{M}^t = s^{\overline{M}}) 
  \;\in\; [0,1]^{|Z|},
\]
donde $\overline{M} = V \setminus M$ indica el complemento del mecanismo.

\begin{definition}[Bipartición MIP]
  Sea $V$ el conjunto de nodos. Una \emph{bipartición} es un par
  $(A, B)$ con $A \cup B = V$, $A \cap B = \emptyset$, $A,B \ne \emptyset$.
  La \emph{MIP} es la bipartición que minimiza:
  \[
    \phi(A,B) = \EMD\bigl(p(Z \mid M),\;p(Z_A \mid M_A) \otimes p(Z_B \mid M_B)\bigr).
  \]
\end{definition}

\begin{definition}[$k$-partición MIP]
  Una $k$-partición de $V$ es una colección $\parti = \{P_1,\ldots,P_k\}$ de
  subconjuntos no vacíos, disjuntos y que cubren $V$.  La $k$-MIP minimiza:
  \[
    \phi(\parti) = \EMD\!\Bigl(p(Z \mid M),\;
                               \bigotimes_{i=1}^k p(Z_{P_i} \mid M_{P_i})\Bigr).
  \]
\end{definition}

\subsection{Earth Mover's Distance analítica}

Para distribuciones de efectos sobre variables binarias condicionalmente
independientes, la EMD admite una forma cerrada~\cite{geomip2025}:
\[
  \EMD(u, v) = \sum_{i=1}^{n} \bigl|u_i - v_i\bigr|,
\]
donde $u_i = P(X_i = 0 \mid \cdot)$ es la probabilidad marginal de que el
nodo $i$ esté \emph{apagado}.  Esta simplificación reduce la complejidad de
$\mathcal{O}(n^3)$ (problema de transporte general) a $\mathcal{O}(n)$.

\subsection{Análisis del espacio de búsqueda}

El número de $k$-particiones de $n$ elementos es el
\emph{número de Stirling de segundo tipo} $S(n,k)$, que crece
aproximadamente como $(k^n)/(k!)$ para $k$ fijo.  La Tabla~\ref{tab:stirling}
muestra su crecimiento:

\begin{table}[H]
  \centering
  \caption{Número de $k$-particiones $S(n,k)$ para distintos $n$ y $k$.}
  \label{tab:stirling}
  \begin{tabular}{r r r r r}
    \toprule
    $n$ & $k=2$     & $k=3$        & $k=4$           & $k=5$            \\
    \midrule
     3  &          3 &            1 &              --  &              --  \\
     5  &         15 &           25 &              10  &               1  \\
    10  &        511 &       9\,330 &          34\,105 &          42\,525 \\
    15  &     16\,383 & 2\,375\,101  & $4.8{\times}10^9$ & $4.2{\times}10^{11}$ \\
    20  & $5{\times}10^5$ & $7.4{\times}10^8$ & $4.5{\times}10^{12}$ & $1.4{\times}10^{16}$ \\
    \bottomrule
  \end{tabular}
\end{table}

La búsqueda exhaustiva es inviable para $n > 6$ incluso para $k=2$.
KGeoMIP explota la estructura del hipercubo para reducir la búsqueda a
$\mathcal{O}(n \cdot 2^n)$ mediante programación dinámica.

% ============================================================
\section{Arquitectura del Software}
\label{sec:arquitectura}
% ============================================================

\subsection{Visión general}

El sistema se organiza en dos módulos principales bajo \texttt{GeoMIP/src/}:

\begin{itemize}
  \item \textbf{Method2\_Dynamic\_Programming\_Reformulation}: núcleo del
        proyecto, contiene las estrategias GeometricSIA, KPartitionSIA,
        QNodes y BruteForce.
  \item \textbf{Method3\_Batch\_Processing}: variante de procesamiento en
        lote sobre todos los TPMs disponibles.
\end{itemize}

\subsection{Diagrama de paquetes}

% TikZ — diagrama de paquetes
\begin{figure}[H]
  \centering
  \begin{tikzpicture}[
    pkg/.style={draw, rounded corners, fill=blue!8,
                minimum width=3.2cm, minimum height=0.7cm, font=\small},
    arr/.style={-{Stealth}, thick, gray}
  ]
    \node[pkg] (manager)    at (0,0)   {\texttt{controllers/manager}};
    \node[pkg] (geo)        at (4,1)   {\texttt{strategies/geometric}};
    \node[pkg] (kpart)      at (4,0)   {\texttt{strategies/kpartition}};
    \node[pkg] (qnodes)     at (4,-1)  {\texttt{strategies/q\_nodes}};
    \node[pkg] (force)      at (4,-2)  {\texttt{strategies/force}};
    \node[pkg] (sia)        at (0,-2)  {\texttt{models/base/sia} (ABC)};
    \node[pkg] (solution)   at (0,-3.5){\texttt{models/core/solution}};
    \node[pkg] (system)     at (4,-3.5){\texttt{models/core/system}};
    \node[pkg] (funcs)      at (0,-5)  {\texttt{funcs/base} (EMD)};

    \draw[arr] (manager) -- (geo);
    \draw[arr] (manager) -- (kpart);
    \draw[arr] (manager) -- (qnodes);
    \draw[arr] (manager) -- (force);
    \draw[arr] (sia) -- (geo);
    \draw[arr] (sia) -- (kpart);
    \draw[arr] (sia) -- (qnodes);
    \draw[arr] (sia) -- (force);
    \draw[arr] (geo)    -- (solution);
    \draw[arr] (kpart)  -- (solution);
    \draw[arr] (qnodes) -- (solution);
    \draw[arr] (force)  -- (solution);
    \draw[arr] (sia)    -- (funcs);
    \draw[arr] (sia)    -- (system);
  \end{tikzpicture}
  \caption{Diagrama de paquetes simplificado de Method2.}
  \label{fig:paquetes}
\end{figure}

\subsection{Diagrama de clases (UML)}

% TikZ — diagrama de clases UML
\begin{figure}[H]
  \centering
  \begin{tikzpicture}[
    cls/.style={draw, fill=white, rounded corners=2pt,
                minimum width=4.5cm, inner sep=4pt, font=\small},
    inh/.style={-{Triangle[open]}, thick},
    dep/.style={-{Stealth}, dashed, gray}
  ]
    \node[cls] (sia) at (0,0) {
      \begin{tabular}{c}
        \textbf{SIA} (ABC)\\
        \hline
        \texttt{+ sia\_preparar\_subsistema()}\\
        \texttt{+ sia\_cargar\_tpm()}\\
        \texttt{\# sia\_gestor: Manager}\\
        \texttt{\# sia\_subsistema: System}\\
        \hline
        \texttt{+ aplicar\_estrategia() [abs]}
      \end{tabular}
    };
    \node[cls] (geo) at (-3.5,-4) {
      \begin{tabular}{c}
        \textbf{GeometricSIA}\\
        \hline
        \texttt{+ tabla\_transiciones}\\
        \texttt{+ memoria\_particiones}\\
        \hline
        \texttt{+ aplicar\_estrategia()}\\
        \texttt{+ calcular\_costos\_nivel()}
      \end{tabular}
    };
    \node[cls] (kp) at (3.5,-4) {
      \begin{tabular}{c}
        \textbf{KPartitionSIA}\\
        \hline
        \texttt{+ k: int}\\
        \hline
        \texttt{+ aplicar\_estrategia()}\\
        \texttt{- \_greedy()}\\
        \texttt{- \_clustering()}
      \end{tabular}
    };
    \node[cls] (sol) at (0,-8) {
      \begin{tabular}{c}
        \textbf{Solution}\\
        \hline
        \texttt{+ perdida: float}\\
        \texttt{+ particion: str}\\
        \texttt{+ estrategia: str}\\
        \texttt{+ n\_nodos: int}\\
        \texttt{+ k: int}\\
        \texttt{+ tiempo\_ms: float @property}
      \end{tabular}
    };

    \draw[inh] (geo) -- (sia);
    \draw[inh] (kp)  -- (sia);
    \draw[dep] (geo) -- (sol) node[midway,left,font=\tiny]{crea};
    \draw[dep] (kp)  -- (sol) node[midway,right,font=\tiny]{crea};
  \end{tikzpicture}
  \caption{Diagrama de clases UML (simplificado).}
  \label{fig:clases}
\end{figure}

\subsection{Patrones de diseño}

\begin{description}[leftmargin=1.5cm]
  \item[\textbf{Template Method}] La clase abstracta \texttt{SIA} define
    \texttt{sia\_preparar\_subsistema()} como operación común; cada subclase
    implementa \texttt{aplicar\_estrategia()}.
  \item[\textbf{Strategy}] Las clases \texttt{GeometricSIA}, \texttt{KPartitionSIA},
    \texttt{QNodes} y \texttt{BruteForce} son estrategias intercambiables que
    comparten la misma interfaz.
  \item[\textbf{Value Object}] \texttt{Solution} encapsula el resultado inmutable
    de la búsqueda, incluyendo los nuevos campos \texttt{n\_nodos}, \texttt{k}
    y la propiedad derivada \texttt{tiempo\_ms}.
\end{description}

% ============================================================
\section{Diseño Algorítmico}
\label{sec:algoritmos}
% ============================================================

\subsection{GeometricSIA (KGeoMIP, $k=2$)}

El algoritmo explota la estructura del hipercubo $\{0,1\}^n$ para calcular
el coste de transición $t(s_0, s')$ de forma incremental, aprovechando que
cada vértice del hipercubo difiere de sus vecinos en exactamente un bit
(operador de volteo).

\begin{algorithm}[H]
  \caption{GeometricSIA --- \texttt{aplicar\_estrategia}}
  \label{alg:geo}
  \begin{algorithmic}[1]
    \Require TPM $\tpm$, estado inicial $s_0$, máscaras cond/alcance/mec
    \Ensure $\parti^* = \arg\min_{(A,B)} \phi(A,B)$, pérdida $\phi^*$
    \State $\text{subsistema} \gets \texttt{sia\_preparar\_subsistema}(\tpm, \text{cond, alcance, mec})$
    \State Inicializar $\texttt{tabla}[s_0 \to s_0] \gets \mathbf{0}^n$
    \For{nivel $\ell = 1, \ldots, n$}
      \For{cada vértice $v$ a distancia $\ell$ de $s_0$}
        \State $v' \gets$ vecino de $v$ más cercano a $s_0$
        \State $\texttt{tabla}[s_0 \to v] \gets \texttt{tabla}[s_0 \to v'] + \Delta(v', v)$
        \Comment{incremento analítico de un bit}
      \EndFor
    \EndFor
    \State $\phi^* \gets +\infty$
    \For{cada bipartición $(A, B)$ de las dimensiones del subsistema}
      \State Calcular $\phi(A,B)$ usando la tabla de costos
      \If{$\phi(A,B) < \phi^*$}
        \State $\phi^* \gets \phi(A,B)$;\; $\parti^* \gets (A,B)$
      \EndIf
    \EndFor
    \State \Return \texttt{Solution}($\phi^*$, $\parti^*$, $k=2$)
  \end{algorithmic}
\end{algorithm}

\subsection{KPartitionSIA (KGeoMIP, $k > 2$)}

\begin{algorithm}[H]
  \caption{KPartitionSIA --- Estrategia greedy encadenada}
  \label{alg:kpart}
  \begin{algorithmic}[1]
    \Require TPM, $s_0$, máscaras, $k$
    \Ensure $\parti^* = \{P_1,\ldots,P_k\}$
    \State Verificar $k \le |V|$; si no, devolver Solution ``Inviable''
    \State Correr \texttt{GeometricSIA} para obtener bipartición inicial $(A_1, A_2)$
    \State $\texttt{grupos} \gets [(A_1, \phi(A_1)), (A_2, \phi(A_2))]$
    \For{$i = 3, \ldots, k$}
      \State $g^* \gets \arg\max_{g \in \texttt{grupos}} \phi(g)$
      \Comment{dividir el grupo con mayor pérdida}
      \State Correr \texttt{GeometricSIA} sobre $g^*$ para obtener $(B_1, B_2)$
      \State $\texttt{grupos} \gets (\texttt{grupos} \setminus \{g^*\}) \cup \{B_1, B_2\}$
    \EndFor
    \State $\phi^* \gets \EMD\bigl(p(Z|M),\;\bigotimes_{g \in \texttt{grupos}} p(Z_g|M_g)\bigr)$
    \State \Return \texttt{Solution}($\phi^*$, grupos, $k$)
  \end{algorithmic}
\end{algorithm}

% ============================================================
\section{Análisis de Complejidad}
\label{sec:complejidad}
% ============================================================

\begin{table}[H]
  \centering
  \caption{Complejidad asintótica de cada estrategia.}
  \label{tab:complejidad}
  \begin{tabular}{l l l l}
    \toprule
    Estrategia & Tiempo & Espacio & Óptima \\
    \midrule
    BruteForce     & $\mathcal{O}(2^n)$         & $\mathcal{O}(n)$       & Sí (bipart.) \\
    QNodes         & $\mathcal{O}(n^2)$          & $\mathcal{O}(n)$       & No           \\
    GeometricSIA   & $\mathcal{O}(n \cdot 2^n)$ & $\mathcal{O}(n \cdot 2^n)$ & Sí (bipart.) \\
    KPartitionSIA  & $\mathcal{O}(k \cdot n \cdot 2^n)$ & $\mathcal{O}(n \cdot 2^n)$ & No (aprox.) \\
    \bottomrule
  \end{tabular}
\end{table}

\subsection{Análisis detallado de GeometricSIA}

La tabla de transiciones tiene $2^n$ entradas, cada una de longitud $n$,
lo que da un espacio $\Theta(n \cdot 2^n)$ bits.  La construcción nivel a
nivel recorre cada vértice del hipercubo una sola vez, dando complejidad
temporal $\Theta(n \cdot 2^n)$.  La posterior búsqueda de la MIP sobre las
biparticiones ocupa $\mathcal{O}(2^{n/2})$ comparaciones gracias a la
simetría del hipercubo~\cite{geomip2025}.

\subsection{Análisis de KPartitionSIA}

KPartitionSIA aplica GeometricSIA exactamente $k-1$ veces sobre sub-sistemas
cada vez más pequeños.  Dado que el subsistema más grande tiene $n$ nodos,
la cota superior es:
\[
  T_k = \sum_{i=0}^{k-2} \mathcal{O}\bigl((n-i) \cdot 2^{n-i}\bigr)
       = \mathcal{O}\bigl(k \cdot n \cdot 2^n\bigr).
\]

% ============================================================
\section{Detalles de Implementación}
\label{sec:implementacion}
% ============================================================

\subsection{Dependencias externas}

\begin{table}[H]
  \centering
  \caption{Dependencias principales del proyecto.}
  \begin{tabular}{l l l}
    \toprule
    Paquete & Versión mín. & Uso \\
    \midrule
    \texttt{numpy}      & 2.0   & Álgebra lineal, arreglos, TPM \\
    \texttt{scipy}      & 1.11  & Clustering jerárquico (\texttt{scipy.cluster}) \\
    \texttt{pyemd}      & 0.6   & EMD exacta (validación) \\
    \texttt{pandas}     & 2.0   & Entrada/salida Excel \\
    \texttt{openpyxl}   & 3.1   & Lectura/escritura de \texttt{.xlsx} \\
    \texttt{colorama}   & 0.4   & Salida en consola con colores \\
    \texttt{pyinstrument} & 4.6 & Profiling (desarrollo) \\
    \texttt{pytest}     & 9.0   & Tests automatizados \\
    \bottomrule
  \end{tabular}
\end{table}

\subsection{Clase \texttt{Solution} (estandarizada)}

Tras la refactorización de la Fase~1, \texttt{Solution} expone los siguientes
campos mínimos requeridos por el benchmark:

\begin{lstlisting}[caption={Campos de Solution tras estandarización.}]
class Solution:
    estrategia: str       # nombre de la estrategia
    perdida: float        # phi = EMD(SO, SP)
    particion: str        # representacion textual de la MIP
    tiempo_ejecucion: float  # tiempo en segundos
    n_nodos: int          # numero de nodos del subsistema (nuevo)
    k: int                # numero de particiones (nuevo)
    distribucion_subsistema: np.ndarray
    distribucion_particion:  np.ndarray

    @property
    def tiempo_ms(self) -> float:
        return self.tiempo_ejecucion * 1000.0
\end{lstlisting}

\subsection{Suite de tests}

Se implementaron 34 tests automáticos en \texttt{tests/}:

\begin{table}[H]
  \centering
  \caption{Resumen de la suite de tests (34 total, 100\% aprobados).}
  \begin{tabular}{l l r}
    \toprule
    Módulo & Cobertura & Tests \\
    \midrule
    \texttt{test\_emd.py}            & Función \texttt{emd\_efecto}           & 7 \\
    \texttt{test\_geometric\_n3.py}  & Tabla~4.2 N3C + ejecución + cross-val  & 10 \\
    \texttt{test\_kpart\_coherence.py} & Campos $k$, $n$, coherencia, inviable & 6 \\
    \texttt{test\_solutions.py}      & Campos mínimos en todas las estrategias & 11 \\
    \midrule
    Total & & 34 \\
    \bottomrule
  \end{tabular}
\end{table}

% ============================================================
\section{Resultados Experimentales}
\label{sec:resultados}
% ============================================================

\subsection{Datasets utilizados}

Los experimentos se ejecutaron sobre 27 archivos CSV de TPMs, con $n$ entre
3 y 22 nodos.  Los archivos con $n \le 20$ se cargan completamente en memoria;
los de $n = 21,22$ se validan sobre una muestra (primeras 1000 filas) por
limitación de RAM.

\subsection{Métricas de evaluación}

\begin{itemize}
  \item \textbf{perdida\_emd}: valor $\phi$ de la partición encontrada.
  \item \textbf{tiempo\_ms}: tiempo de CPU en milisegundos.
  \item \textbf{memoria\_mb}: pico de memoria RSS en megabytes.
  \item \textbf{convergio}: booleano; falso si el caso excedió el timeout o
        fue declarado inviable.
\end{itemize}

\subsection{Tablas de resultados}
\label{sec:tablas}

% Tablas generadas automáticamente por docs/gen_tables.py
% (se insertan durante la compilación final)
% Tabla pendiente: ejecutar docs/gen_tables.py despues del benchmark
% python docs/gen_tables.py

   % generado por gen_tables.py

\begin{table}[H]
  \centering
  \caption{Resultados representativos para $n=10$, $k=2$.
           Estado inicial: \texttt{1000000000}.
           (Tabla de ejemplo; resultados reales en \texttt{benchmark\_*.xlsx}).}
  \label{tab:n10k2}
  \begin{tabular}{l l r r r r}
    \toprule
    Subsistema (purview/mec) & Estrategia & $\phi$ & $t$ (ms) & Mem (MB) & Convergió \\
    \midrule
    ABCDEFGHIJ / ABCDEFGHIJ & GeometricSIA & --- & --- & --- & --- \\
    ABCDEFGHIJ / ABCDEFGHIJ & QNodes       & --- & --- & --- & --- \\
    \multicolumn{6}{c}{\emph{(datos del archivo benchmark generado)}} \\
    \bottomrule
  \end{tabular}
\end{table}

\subsection{Validación de correctitud}

GeometricSIA se validó contra BruteForce para el sistema N3A (3 nodos):
la diferencia en $\phi$ fue $< 10^{-6}$ en todos los estados iniciales
probados.  Además, los ocho costos de transición de la Tabla~4.2 del
artículo GeoMIP~\cite{geomip2025} se reproducen con error $< 10^{-4}$
(ver tests en \texttt{test\_geometric\_n3.py}).

% ============================================================
\section{Limitaciones y Trabajo Futuro}
\label{sec:limitaciones}
% ============================================================

\subsection{Limitaciones conocidas}

\begin{enumerate}
  \item La heurística greedy de KPartitionSIA no garantiza monotonicidad
        $\phi(k+1) \le \phi(k)$ en todos los casos.  En los experimentos
        sobre N3--N5 se observó esta propiedad en solo $\approx 33\,\%$
        de los casos.
  \item Para $n > 15$, la tabla de transiciones ocupa más de 4 GB,
        superando la RAM de máquinas típicas.
  \item Los archivos $N21$/$N22$ solo se validan parcialmente (1000 filas);
        el benchmark completo requiere paralelismo o muestreo de estados.
\end{enumerate}

\subsection{Trabajo futuro}

\begin{itemize}
  \item Implementar memoización de filas de la tabla de costos usando
        tablas hash para reducir el espacio a $\mathcal{O}(n^2)$.
  \item Paralelizar el cálculo de la tabla con \texttt{ProcessPoolExecutor}
        o GPU (POT/CuPy) según el addendum~\ref{app:gpu}.
  \item Extender KQNodes para que aplique $k$ cortes greedy de forma
        coordinada (actualmente solo aplica en GeoMIP).
  \item Generar y validar los TPMs de $n=23,24,25$ usando
        \texttt{generate\_large\_tpms.py}.
\end{itemize}

% ============================================================
\appendix
% ============================================================

\section{Demostración: EMD analítica para efectos binarios}
\label{app:emd}

\begin{proposition}
  Sea $X = (X_1,\ldots,X_n)$ con $X_i \in \{0,1\}$.  Si $X_1,\ldots,X_n$
  son condicionalmente independientes bajo $p$ y $q$, entonces con la métrica
  de Hamming como suelo:
  \[
    \EMD(p, q) = \sum_{i=1}^{n} |p(X_i = 0) - q(X_i = 0)|.
  \]
\end{proposition}

\begin{proof}[Demostración (esquema)]
  La independencia condicional factoriza las distribuciones como
  $p(x) = \prod_i p_i(x_i)$ y $q(x) = \prod_i q_i(x_i)$.
  El problema de transporte óptimo se descompone en $n$ subproblemas
  unidimensionales independientes, cada uno con coste $|p_i(0) - q_i(0)|$
  (la única masa que debe moverse entre los dos estados posibles del nodo $i$).
  La suma da el resultado.
\end{proof}

\section{Plan de paralelización GPU (Fase 4B)}
\label{app:gpu}

El documento \texttt{ADDENDUM\_GPU\_PARALLELIZACION.md} especifica la
implementación opcional de:

\begin{itemize}
  \item \textbf{CPU multicore}: \texttt{ProcessPoolExecutor} para paralelizar
        el cálculo de la tabla de transiciones por nivel del hipercubo.
  \item \textbf{GPU EMD}: librería POT (\texttt{ot.emd}) con backend CUDA
        para lotes de distribuciones.
  \item \textbf{GPU Hamming}: CuPy para distancias de Hamming vectorizadas.
  \item \textbf{GPU clustering}: cuML \texttt{AgglomerativeClustering} para
        KPartitionSIA con $n > 12$.
\end{itemize}

Estas optimizaciones \textbf{solo se implementan} si el benchmark de la
Fase~3 justifica el esfuerzo (tiempos $> 60$ s o errores de memoria para
$n \ge 15$).

\section{Uso de Inteligencia Artificial Generativa}
\label{app:ia}

Durante el desarrollo del proyecto se utilizó \textbf{Cursor AI} (Claude
Sonnet 4.6) como asistente de codificación.  Las áreas de uso fueron:

\begin{description}
  \item[Diagnóstico y corrección de bugs:] Identificación del bug en
    \texttt{force.py} y \texttt{q\_nodes.py} (argumento \texttt{tpm}
    faltante), y la deprecación de \texttt{np.infty} en NumPy~2.0.
  \item[Generación de tests:] Esqueleto inicial de los 34 tests;
    todos los valores esperados fueron verificados manualmente contra
    el Cuadro~4.2 del artículo GeoMIP y los ejemplos del documento
    \texttt{Ejemplos.md}.
  \item[Documentación:] Estructura del presente manual.
  \item[Benchmark:] Estructura de \texttt{benchmark.py} y el validador
    \texttt{validate\_tpms.py}.
\end{description}

Todas las decisiones algorítmicas y la verificación de correctitud fueron
realizadas por el equipo.

% ── Bibliografía ─────────────────────────────────────────────────────────────
\bibliographystyle{plain}
\begin{thebibliography}{9}

\bibitem{tononi2016}
G.~Tononi, M.~Boly, M.~Massimini, C.~Koch,
\textit{Integrated information theory: from consciousness to its physical
substrate},
Nature Reviews Neuroscience, 17(7):450--461, 2016.

\bibitem{geomip2025}
Autor(es) del equipo,
\textit{GeoMIP: A Geometric Approach to the Minimum Information Partition
Problem},
Proyecto de grado, Universidad de Caldas, 2025.

\bibitem{tononi2004}
G.~Tononi,
\textit{An information integration theory of consciousness},
BMC Neuroscience, 5:42, 2004.

\bibitem{emd1998}
Y.~Rubner, C.~Tomasi, L.J.~Guibas,
\textit{A metric for distributions with applications to image databases},
ICCV 1998.

\end{thebibliography}

\end{document}
